\documentclass{article}
\usepackage[utf8]{inputenc}
\usepackage[T1]{fontenc}
\usepackage{geometry}
\geometry{
    a4paper,
    left=2.5cm,
    right=2.5cm,
    top=2.5cm,
    bottom=2.5cm
}
\usepackage{graphicx}
\usepackage{float}
\usepackage{hyperref}
\usepackage[english]{babel}
\usepackage{amsmath}
\usepackage{amssymb}
\usepackage{array}
\usepackage{booktabs}
\usepackage{caption}
\usepackage{subcaption}
\usepackage{enumitem}
\usepackage{fancyhdr}
\usepackage{titlesec}
\usepackage{url}
\usepackage{textcomp}
\usepackage{xcolor}
\usepackage{listings}
\usepackage{multirow}
\usepackage{longtable}
\usepackage{tabularx}
\usepackage{times}

% Set up page style
\pagestyle{fancy}
\fancyhf{}
\rfoot{\thepage}
\renewcommand{\headrulewidth}{0pt}

% Title formatting
\titleformat{\chapter}[display]{\normalfont\bfseries}{\chaptertitle\thechapter}{0pt}{\Huge}
\titlespacing*{\chapter}{0pt}{50pt}{40pt}

% Document info
\title{Template Document}
\author{Academic Support Agent}
\date{\today}

\begin{document}
\maketitle

%%=============================================================================
%% Capa e folha de rosto
%%=============================================================================
\maketitle
%%=============================================================================
%% Cataloga o (obrigat rio para mestrado) e Folha de aprova o
%%=============================================================================
%Somente obrigat rio para disserta o, para TG, remover as linhas	77	%
%Como a CIP vai ser impressa atr s da p gina de rosto, as margens inner e outer	
%devem ser invertidas.

\newgeometry{inner=20mm,outer=30mm,top=30mm,bottom=20mm}	
\makeCIP{autor@email} %email do autor		
\restoregeometry

%Se for usar a cataloga o gerada pelo gerador do site da biblioteca comentar as linhas
%\includeCIP{CIP.pdf}
	
%%=============================================================================
%folha de aprova o
%%=============================================================================
\makeapprove

%%=============================================================================
%% Dedicat ria (opcional)
%%=============================================================================
% Content from pessoal/dedicatoria would be included here

%%=============================================================================
%% Agradecimentos (opcional)
%%=============================================================================
% Content from pessoal/agradecimento would be included here

%%=============================================================================
%% Ep grafe (opcional)
%%=============================================================================
% Content from pessoal/epigrafe would be included here

%%=============================================================================
%% Resumo
%%=============================================================================
\begin{abstract}
		
Aqui voc  escreve o resumo. Lembrando no m ximo 250 palavras para tcc e 500 palavras para tese ou disserta o.

\end{abstract}

%%=============================================================================
%% Abstract
%%=============================================================================
% resumo na outra l ngua
% como parametros devem ser passados o titulo, o nome do curso,
% as palavras-chave na outra l ngua, separadas por v rgulas, o m s em ingl s
%o a sigla do dia em ingl s: st, nd, th ...
\begin{englishabstract}
{Abstract title}
%%============ MANTER ESSA LINHA ============================================
\ %%============ MANTER ESSA LINHA ==========================================
%%============ MANTER ESSA LINHA ============================================
\ %%============ MANTER ESSA LINHA ==========================================
%%============ MANTER ESSA LINHA ============================================
\ %%============ MANTER ESSA LINHA e deixar uma em branco ===================


Here you write the summary. Remembering a maximum of 250 words for tcc and 500 words for thesis or dissertation.

\end{englishabstract}

%% Lista de Ilustra es (opc)
%% Lista de S mbolos (opc)
%% Lista de Anexos e Ap ndices (opc)

%%=============================================================================
%% Lista de figuras (comentar se n o houver)
%%=============================================================================

%%=============================================================================
%% Lista de tabelas (comentar se n o houver)
%%=============================================================================


%%=============================================================================
%% Lista de Ap ndices (comentar se n o houver)
%%=============================================================================
%\listofappendix

%%=============================================================================
%% Lista de Anexos (comentar se n o houver)
%%=============================================================================
%\listofannex

%%=============================================================================
%% Lista de abreviaturas e siglas
%%=============================================================================
 %o parametro deve ser a abreviatura mais longa

% Content from siglas/lista would be included here


%%=============================================================================
%% Lista de simbolos (opcional)
%%=============================================================================
%Simbolos devem aparecer conforme a ordem em que aparecem no texto
% o parametro deve ser o s mbolo mais longo
%\begin{listofsymbols}{teste}
%  \item [$\varnothing$] vazio
 % \item [$\Gamma$]  Gama
%  \item [$\forall$] Para todo
%\end{listofsymbols}

%%=============================================================================
%% Sum rio
%%=============================================================================












%%=============================================================================
%% In cio da disserta o
%%=============================================================================
\setlength{\baselineskip}{1.5\baselineskip}

%Adiciona cada capitulo
% Content from capitulos/introducao would be included here

% Content from capitulos/desenvolvimento would be included here

% Content from capitulos/conclusao would be included here


\setlength{\baselineskip}{\baselineskip}

%%=============================================================================
%% Refer ncias
%%=============================================================================
\begin{flushleft}


\end{flushleft}

%IMPORTANTE: Se precisar usar alguma se o ou subse o dentro dos ap ndices ou
%anexos, utilizar o comando \tocless para n o adicionar no Sum rio
%Exemplos: 
% \tocless\section{Hist rico}
%%=============================================================================
%% Ap ndices
%%=============================================================================


% Content from apendices/apendicea would be included here

%% Content from apendices/apendiceb would be included here


%%=============================================================================
%% Anexos
%%=============================================================================
\annex
% Content from anexos/anexoa would be included here
\end{document}